\documentclass[oneside,12pt]{book}
    % The option oneside is required for UNLV theses.
    % The 12pt is optional -- UNLV will allow 10pt too.
    % 10pt is the default, so to switch to 10pt, just delete 12pt in the above.
\usepackage{amsthm,amssymb,amsmath,graphicx} % Optional packages.
\usepackage{setspace,UNLVthesis} % Required packages.  Look at UNLVthesis.sty to see how
    % LaTeX is instructed to set things up.  This file may need some tweaking.
    % setspace.sty is not normally part of MikTeX.  It can be obtained from www.ctan.org.
    % Do a search on setspace.sty.
    % The files setspace.sty and UNLVthesis.sty should be in the same directory
    % as this file (or in a directory of MikTeX where LaTeX will know to find it
    % -- for example, where other style files are).
\pagestyle{unlv} %  This is defined in UNLVthesis.sty.  Headings are empty except for page numbers,
    % and the page numbers are the same size as the text.  Most documents use a different size so that
    % it it is difficult to mistake it as part of the text.

%These define the format and numbering of theorem like environments.
\newtheorem{theorem}{Theorem}
\newtheorem{corollary}[theorem]{Corollary}%Corollaries and Lemmas are numbered as theorems.
\newtheorem{lemma}[theorem]{Lemma}

%These define the format and numbering of definition like environments.
\theoremstyle{definition}%This environment is not in italics, like theorems are.
\newtheorem{definition}{Definition}
\newtheorem*{introduction}{Introduction}%The * means it is unnumbered.
\newtheorem*{conclusion}{Conclusion}

% Put definitions here.  For example, suppose you often use the Greek characters
% alpha, beta, etc., which in LaTeX are \alpha, \beta, etc. (in math mode only).
% Then it may be easier to create shortcuts for these commands, such as:
%\def\aa{\alpha}
%\def\bb{\beta}
% Now, instead of typing \alpha, we can type \aa.
% Here's another I often use:
\def\Bbb#1{{\mathbb #1}} % \Bbb is an obsolete command, but I'm old
% and still use it, so I define it to do what it used to do.  The usage
% is like \Bbb R, which will produce a blackboard bold R, and is literally
% translated to {\mathbb R}.  Note that this command includes a single
% argument.

\begin{document}

\pagenumbering{roman}% Items before Chapter one have roman numbers (if any).
\thispagestyle{empty}%This page has no page number.
 \begin{center}
 THE UNLV THESIS STYLE PACKAGE\\*[48pt]% The title can be no more than 80 characters -- UNLV rule.
 %  The command \\*[48pt] means carriage return with a 48pt gap.  These numbers can be adjusted
 %  to improve the look.

\normalsize by\\*[48pt]

 Arthur Baragar\\*[48pt]

 Bachelor of Science\\*[-12pt]%Single space ...
 University of Alberta, Canada\\*[-12pt]
 1986\\*[48pt]


 A thesis submitted in partial fulfillment\\*[-12pt]
 of the requirements for the\\*[48pt]


 {\bf Master of Science Degree} \\*[-12pt]
 {\bf Department of Mathematical Sciences}\\*[-12pt]
 {\bf College of Sciences}\\*[36pt]

\normalsize
 {\bf Graduate College}\\*[-12pt]
 {\bf University of Nevada, Las Vegas}\\*[-12pt]
 {\bf May 2002}
\end{center}
 %Replace this file name with the name of your title page.
    % A copyright statement is optional and would be placed here.
    % The copyright page has no page number -- the title page is always page i and the
    % Thesis Approval page is always page ii.
\newpage \setcounter{page}{3} %The Thesis Approval page is page ii.  It is inserted separately.
\chapter*{ABSTRACT}
\addcontentsline{toc}{schapter}{ABSTRACT}% This command adds the Abstract to the table of contents.

\begin{center}
\textbf{THE UNLV THESIS STYLE PACKAGE}

by

Arthur Baragar

 $\langle$Name$\rangle$, Examination Committee
 Chair\\*[-12pt]%Single spaced.
 Professor of Mathematical Sciences \\*[-12pt]
 University of Nevada, Las Vegas  %
 \end{center}

The abstract would be here.  The purpose of this package is to
give some advice about using \LaTeX to prepare their thesis, and
to illustrate, via an example (this document), what the package
produces, as well as provide a template.  This package should be
used in conjunction with the source files: UNLVthesisTemplate.tex,
SampleTitlepage.tex, SamplePreface.tex, SampleAbstract.tex,
SampleChapter1.tex, SampleChapter2.tex, SampleChapter2.tex,
SampleBibliography.tex, and SampleVita.tex.
 % This is where the abstract goes.

\tableofcontents %inserts table of contents
\listoffigures \addcontentsline{toc}{schapter}{\listfigurename}% Comment these out if there are no figures or tables.
\listoftables \addcontentsline{toc}{schapter}{\listtablename}%

\chapter*{ACKNOWLEDGEMENTS}
\addcontentsline{toc}{schapter}{ACKNOWLEDGEMENTS}

I would like to thank Doina Bein for letting me see what she did
when she wrote her thesis.  I would like to thank Kensaku Umeda,
who helped me with the UNLV requirements.
  %Acknowledgements come after the tables.

\newpage % Do not remove this command.  It's there to make sure the page numbering is correct.
\pagenumbering{arabic} %Chapter 1 begins on page 1.

\chapter{INTRODUCTION}% Chapter titles are capitalized -- UNLV rule.

Writing a graduate thesis is very time consuming.  One of the
obstacles placed in front of the student is the set of
requirements for the look of thesis.  These requirements look like
they date from the middle ages, but there's nothing you or I can
do about that.  Writing a thesis in mathematics or some other
discipline that requires mathematical symbols is even more
difficult, because most typesetting programs do not deal with them
very well.  One program that was designed to handle mathematical
symbols is \LaTeX.  This program is a publishing grade program.
For example, I have published a book that was typeset using \LaTeX
\cite{B1} and almost all journals that I submit to request that
source files be in \LaTeX.

Unfortunately, \LaTeX\ is much more sophisticated than what the
average graduate student requires, and none of its standard
packages look anything like the required UNLV format (though every
one of them looks better). On the other hand, much of this can be
automated by a \TeX\! pert. That's what this package is about.  It
describes how to use the UNLVthesis.sty style file to format
theses into the UNLV format. Unfortunately, I do not make the
rules, so there may be things I haven't anticipated or dealt with
correctly, so some tweaking may be required.  I welcome questions,
though I may not deal with them promptly.  I can be reached at
baragar@unlv.edu.

There are some requirements that are very difficult to adhere to.
For example, the font style is to be ``a standard ... computer
font, such as Bookman, Courier, Arial(Helvetica), or Times.''
Though it is possible to change the font in \LaTeX\ documents, I
prefer to not do so.  The \LaTeX\ font is Computer Modern Roman,
and very few would notice that it is not Times.
 %Replace these with your chapters.
\chapter{WHERE DO I GET \LaTeX?}

\LaTeX\ is a freeware application.  It was originally designed by
Donald Knuth at Stanford, and is being maintained by the (mostly
European) academic community.  There are versions that run on PC,
Mac, or unix systems.  Though the compilers are platform
dependent, the source files are not.

\section{Obtaining a PC version of \LaTeX}

The PC version of \LaTeX\ is MikTeX.  It can be obtained from
\[
\hbox{www.miktex.org}
\]
As mentioned, it is free.  I strongly recommend the use of the
`front end' editor WinEdt, which is shareware and can be obtained
from
\[
\hbox{www.winedt.com}
\]
The cost is nominal (currently \$40).  Their website also has
links to the necessary downloads.  Besides MikTeX and WinEdt, one
should get Ghostscript and Ghostview (go to the WinEdt website for
links).  These too are free.

\section{Macintosh versions}

Sorry, you're on your own.

\section{Unix versions}

Again, you're on your own.  However, chances are you're using unix
on a mainframe, in which case someone else has probably already
installed \LaTeX.

\section{Getting started}

I suggest using the electronic versions of these files as a
template.  Begin by editing it.  But keep a copy of these, since
in the next chapter, I do several things that are probably
representative of the sort of things you probably will want to do.
You can look at the .tex file to see how I did them.

\section{Resources}

There are several books that I recommend.  I like {\it The \LaTeX\
Companion} by Gossens, Mittelbach, and Samarin \cite{G}, but this
is probably more for an expert. I used to prefer {\it A document
preparation system \LaTeX} by Leslie Lamport \cite{L}, which means
it is more suited to the beginner. I've also heard good things
about {\it A guide to \LaTeX2e, Document Preparation for Beginners
and Advanced Users}, by Helmut Kopka and Patrick Daly \cite{K-D},
but I think this too is for the expert. Finally, one can look at
the starter page for the website
\[
\hbox{www.ctan.org}
\]
This includes links to {\it The (Not so) short introduction to
\LaTeX2$\epsilon$}, by Tobias Oiteker, et al, a 110 page free
document \cite{O}. I haven't tried using this document, but it
looks pretty good.

\chapter[A sample document, eh?]{A SAMPLE DOCUMENT}%Titles must be capitalized.

Before we begin, take a quick look at how this chapter appears in
the table of contents.  This is because of the `optional' argument
in the chapter entry.  Also, while reading this chapter, please
also read its source file SampleChapter3.tex.  Many of the
comments refer to what is written in the source file.

A double return is interpreted as a new paragraph. A single return
is ignored (look at the source file).

To type an equation in the line, put it between dollar signs.  For
example, $a^2+b^2=c^2$.  If it's very important, display the
equation by using backslash square braces:
\[
E=mc^2.
\]
If the equation is to be referenced, then do the following:
\begin{equation}\label{Equation 1}%This is the label of the equation -- LaTeX will number it.
a^n+b^n=c^n.
\end{equation}
(That instruction doesn't make much sense unless you are also
looking at the .tex file, which in this case, is
SampleChapter3.tex.)  Equation (\ref{Equation 1}) has no
non-trivial solutions in the integers. Some results are important:

\begin{theorem}[Fermat's Last Theorem]\label{theorem 1}  The equation
\[
x^n+y^n=z^n
\]
has no solutions $x,y,z\in \Bbb Z$ for $n\geq 3$ and $xyz\neq 0$.
\end{theorem}

Theorem \ref{theorem 1} was proved by Wiles.  The proof does not
appear in \cite{B1}.

Sometimes you might want to have aligned equations, like the
following:
\begin{align*}%the * means this equation doesn't have a label.
\hbox{Force}&=(\hbox{mass})(\hbox{acceleration}) \\ %the \\ tells it to insert a line break.
F&=ma. % the & is the alignment character.
\end{align*}
The * in \verb \begin{align*} \ in the .tex file means that these
equations have no numbers.  If they are to be numbered, leave the
* out:
\begin{align}
\hbox{Force}&=(\hbox{mass})(\hbox{acceleration}) \\
F&=ma.
\end{align}
If you want only one of these equations numbered, try the
following:
\begin{align}
\hbox{Force}&=(\hbox{mass})(\hbox{acceleration}) \notag\\ %the \notag means no tag here.
F&=ma.
\end{align}

A picture is shown in Figure \ref{fig1}.
\begin{figure}[h]%This [h] tells LaTeX to try to put the picture here.
                 %Without it, it will go to the top of the next page.
\begin{center}
{\mbox{\includegraphics[height=140pt]{fig1.eps}}}
\end{center}
\caption{\label{fig1}Isn't this a nice picture?}
\end{figure}% Be careful with empty lines -- remember, they mean new paragraph.
    % There are no empty lines before and after the figure environment since
    % I do not want a new paragraph here.
For more information about graphics, consult documentation on
graphicx (run a search on the www.ctan.org site).  Support of
graphics is really a function of the dvi viewer.  As far as I
know, the only type of graphic supported by the standard viewer is
the .eps format.  Adobe Illustrator\footnote{Adobe Illustrator is
a rather expensive product.  This is the way one does footnotes.
More text here. } can convert most graphics files to .eps files.
Figure \ref{fig1}, though, was produced using Maple and saved as a
\LaTeX\ file (which includes a .eps component).

Though I always use Illustrator\footnote{Another footnote.} to
create pictures, I understand that any printable graphic file can
be converted to .eps format. The idea is this -- if it is
printable, then it can be printed on a postscript printer.  Any
computer can be set up to print to a postscript printer, even if
not connected to one.  So, obtain a generic postscript printer
driver (free from Adobe at www.adobe.com) and set it up to print
to a file.  While setting it up, choose the option to send files
to the ``printer'' in .eps format (the default is .ps format).  I
think I saw full instructions for this in \cite{O}.

For those interested, the commands in Maple that produces
\ref{fig1} are:
\begin{trivlist}
\item plot(\{[2*cos(t),sin(t), t=0..2*Pi],
[-1+cos(t),.2+sin(t),t=0..2*Pi]\}, axes=none, scaling=constrained,
color=black);
\end{trivlist}

\begin{thebibliography}{99}
 \addcontentsline{toc}{schapter}{\bibname}% This adds the Bibliography to the table of contents.
 \begin{spacing}{1} %This makes the bibliography single spaced.
\bibitem{B1}
    A. Baragar, {\it A Survey of Classical and Modern
    Geometries: with Computer Activities}, Prentice Hall, Upper Saddle
    River, NJ, 2001.

\bibitem{B2}
    A. Baragar, Rational points on K3 surfaces in $\Bbb P^1\times
    \Bbb P^1\times \Bbb P^1$, {\it Math. Ann.} {\bf 305}, 541 -- 558
    (1996).  % This is an example of how to reference an article.

\bibitem{G}
    Goossens, Michel, Frank Mittelbach, Alexander Samarin,
    {\it The \LaTeX Companion},
    Addison-Wesley Publishing Co., Reading, MA (1994).

\bibitem{K-D}
    Kopka, Helmut, Patrick Daly, {\it A Guide to \LaTeX2e,
    Document Preparation for Beginners and Advanced Users}, 2nd
    Edition, Addison-Wesley Publishing Co., Reading, MA (1995).

\bibitem{L}
    Lamport, Leslie,
    {\it A Document Preparation System \LaTeX, User's guide and reference manual},
    2nd edition, Addison-Wesley Publishing Co., Reading, MA (1994).

\bibitem{O}
    Oetiker, Tobias, Hubert Partl, Irene Hyna, Elisabeth Schlegl,
    {\it The Not So Short Introduction to \LaTeX2e}, (08/09/01),
    $<$http://www.ctan.org/tex-archive/info/lshort/english/lshort.pdf$>$,
    (04/16/02). % This is how one references a web document.  The first
    % date is the date it was "published," while the second date is the
    % date it was accessed.
 \end{spacing}
\end{thebibliography}
 %There are some special commands in this file.
\begin{spacing}{1} %This makes this page single spaced.
\thispagestyle{plain} % This puts the page number at the bottom.
\begin{center}
 VITA\addcontentsline{toc}{schapter}{VITA}\\*[2\baselineskip]% Don't use \chapter* -- the Vita begins
    % at the top of the page.
 Graduate College \\
 University of Nevada, Las Vegas \\*[\baselineskip]
 Arthur Baragar\\*[\baselineskip]
\end{center}

%\noindent Local Address:\\ % Include this if different from your home address.
%***\\*[2\baselineskip]

\noindent Home Address:

 8815 Wallaby Lane

 Las Vegas, Nevada 89123\\*[.5\baselineskip]

\noindent Degrees:

 Bachelor of Science, Honors Mathematics, 1985

 University of Alberta, Edmonton, Canada\\*[.5\baselineskip]

 Doctor of Philosophy, Mathematics, 1991

 Brown University, Providence, Rhode Island\\*[.5\baselineskip]

%\noindent Special Honors and Awards: % Omit, if none.

%\noindent Publications:  % Omit, if none.

\noindent  Dissertation Title: %(Or Thesis title)
The Markoff Equation and Equations of Hurwitz\\*[.5\baselineskip]

\noindent  Dissertation Examination Committee: %(Or Thesis ...)

 Chairperson, Dr.~Joseph H. Silverman, Ph.D. %Put ~ after a period if the period should not
    % be interpreted as an end of sentence.  This will restrict how much space can be there,
    % and will prevent a line break there.  This is not necessary if there is no space after
    % the period, for example, in www.unlv.edu

 Committee Member, % I don't remember anymore ...

 Committee Member,

 Graduate Faculty Representative,

\end{spacing}
  %The Vita is the last page.

\end{document}
